\documentclass[]{resume}

\usepackage{XCharter}
\usepackage{url}
\usepackage[hidelinks]{hyperref}
\usepackage{xcolor}
\usepackage{fontawesome5}
\usepackage{geometry}
\usepackage{enumitem}
\usepackage{array}

% ---------- Layout ----------
\geometry{left=0.34in,right=0.34in,top=0.46in,bottom=0.34in}
\setlength{\tabcolsep}{4pt}
\setlength{\parskip}{0pt}
\setlist[itemize]{leftmargin=1.2em, itemsep=1pt, topsep=2pt, parsep=0pt, partopsep=0pt}

% resume.cls auto-prints \name/\address at \begin{document}; disable it since we use a custom header.
\makeatletter
\renewcommand{\printname}{}
\renewcommand{\printaddress}[1]{}
\makeatother
\newcommand{\sectionvspace}{\vspace{1.5pt}}

% ---------- Header ----------
\definecolor{HeaderOrange}{HTML}{F36C21}
\definecolor{HeaderLight}{HTML}{F7F7F7}

\newcommand{\MakeHeader}[8]{%
\noindent
{\centering \bfseries \LARGE #1 \par}
\vspace{2pt}

\noindent
\colorbox{HeaderOrange}{%
  \parbox{\textwidth}{%
    \vspace{2pt}
    \centering
    {\color{white}\bfseries\normalsize #2}
    \vspace{2pt}
  }%
}\par

\noindent
\colorbox{HeaderLight}{%
  \parbox{\textwidth}{%
    \vspace{2pt}
    \small
    \faPhone\ #3 \hfill
    \faEnvelope\ \href{mailto:#4}{#4} \hfill
    \faMapMarker*\ #5 \hfill
    \faLinkedin\ \href{#6}{kexinchu} \hfill
    \faGithub\ \href{#7}{kexinchu} \hfill
    \faGraduationCap\ \href{#8}{Google Scholar}
    \vspace{2pt}
  }%
}\par
}

\begin{document}

\MakeHeader
{Kexin Chu}
{Ph.D. Researcher in ML Systems -- University of Connecticut}
{(+1) 959-995-3930}
{kexin.chu@uconn.edu}
{Farmington, CT}
{https://www.linkedin.com/in/kexin-chu-69a481192}
{https://kexinchu.github.io/}
{https://scholar.google.com/citations?user=ZIdS3d0AAAAJ}

% ===================== SUMMARY =====================
\begin{rSection}{Summary}\sectionvspace
\begin{itemize}[leftmargin=1.2em, itemsep=0.15em, topsep=0.2em]
    \item Ph.D. student in Computer Science at UConn, focused on ML systems for LLM/MoE serving, KV-cache management, high-performance ANNS, and inference correctness.
    \item Research spans LLM serving efficiency, multimodal retrieval, and deterministic inference -- with results published at and under review at top venues (ICDCS, SOSP).
\end{itemize}
\end{rSection}

% ===================== EDUCATION =====================
\begin{rSection}{Education}\sectionvspace
\textbf{University of Connecticut}, CT, USA  (GPA:3.91) \hfill \textit{Aug 2024 -- Present}\\
Ph.D. in Computer Science and Engineering

\textbf{Hefei University of Technology}, China \hfill \textit{Aug 2017 -- Jun 2020}\\
M.S. in Electrical Engineering

\textbf{Hefei University of Technology}, China \hfill \textit{Aug 2013 -- Jun 2017}\\
B.S. in Electrical Engineering
\end{rSection}

% ===================== PUBLICATIONS =====================
\begin{rSection}{Publications}\sectionvspace
\item K. Chu, et al. \textit{Dynamic OOD Graph ANNS for Multimodal RAG: Interference-Aware Continuous Repair}. Under review at SOSP 2026.
\item K. Chu, et al. \textit{MCaM: Efficient LLM Inference with Multi-tier KV Cache Management}. ICDCS 2025.
\item K. Chu, et al. \textit{SafeKV: System-Level Co-Design of Privacy Enforcement and KV-Cache Management for LLM Serving}. Under review at EuroSys.
\item K. Chu, et al. \textit{FlexMoE: Adaptive Expert Prefetching and Cache-Aware Routing for Fast MoE Inference}. Under review at ICCAD 2026.
\item K. Chu, et al. \textit{DynaExq: Runtime-Aware Expert Quantization for Efficient MoE Inference}. Under review at ICCAD 2026.
\end{rSection}

% ===================== RESEARCH =====================
\begin{rSection}{Research Projects (UConn)}\sectionvspace

\textbf{Deterministic LLM Inference} \hfill \textit{Mar 2026 -- Present}
\begin{itemize}[leftmargin=1.2em, itemsep=0.2em, topsep=0.2em]
    \item Investigating sources of non-determinism in LLM inference across GPU parallelism, floating-point reductions, and runtime scheduling that cause identical inputs to produce different outputs across runs.
    \item Designing a system to enforce bit-exact reproducibility for LLM inference without sacrificing throughput, targeting production serving systems (vLLM/SGLang) and safety-critical applications.
\end{itemize}

\vspace{0.5ex}
\textbf{Dynamic OOD Graph ANNS for Multimodal RAG} \hfill \textit{Nov 2025 -- Apr 2026}
\begin{itemize}[leftmargin=1.2em, itemsep=0.2em, topsep=0.2em]
    \item Identified that cross-modal (text-over-image/video) OOD queries require up to 9.2$\times$ more graph node visits than in-distribution queries, with ongoing dynamic updates continuously degrading the reachability structure.
    \item Designed an interference-aware continuous repair system with online repair-worthiness prediction, grouped local repair, and a latency-aware controller that bounds tail-latency impact during live serving.
    \item Stabilized recall at 93.7\%--97.8\% across three cross-modal datasets over 5-hour dynamic workloads; P99 latency spikes reduced by over 2$\times$ vs.\ periodic-repair baselines. \textit{Under review at SOSP 2026.}
\end{itemize}

\vspace{0.5ex}
\textbf{SafeKV: Privacy-Aware KV-Cache Management for Shared LLM Serving} \hfill \textit{Apr 2025 -- Aug 2025}
\begin{itemize}[leftmargin=1.2em, itemsep=0.2em, topsep=0.2em]
    \item Designed a cache visibility control mechanism for multi-tenant KV reuse in shared LLM serving.
    \item Combined asynchronous privacy enforcement with metadata-based gating for cache promotion and reuse.
    \item Reduced timing-based KV side-channel attack success by 94\% while improving throughput by 2.66$\times$ versus strict isolation. \textit{Under review at EuroSys.}
\end{itemize}

\vspace{0.5ex}
\textbf{DynaExq: Runtime-Aware Expert Quantization for MoE Inference} \hfill \textit{Jan 2025 -- Apr 2026}
\begin{itemize}[leftmargin=1.2em, itemsep=0.2em, topsep=0.2em]
    \item Identified that static expert quantization misallocates precision under a fixed HBM budget due to heavy-tailed, task-dependent activation; designed a runtime precision residency manager that promotes/demotes expert precision based on observed router traffic.
    \item Built non-blocking stream-parallel precision migration with pool-based memory management and admission control, keeping peak GPU usage within a hard 48\,GB HBM cap.
    \item Recovered 79\% of INT2-to-INT4 accuracy gap (73.58\%$\to$77.15\% on Qwen3-MoE-80B) and achieved 2.73$\times$ throughput vs.\ offloading baselines. \textit{Under review at ICCAD 2026.}
\end{itemize}

\vspace{0.5ex}
\textbf{FlexMoE: Adaptive Expert Prefetching for MoE Inference} \hfill \textit{Aug 2024 -- Apr 2026}
\begin{itemize}[leftmargin=1.2em, itemsep=0.2em, topsep=0.2em]
    \item Identified a U-shaped overlap--pollution tradeoff in fixed-horizon cross-layer prefetching that inflates end-to-end latency by up to 2.5$\times$ across batch sizes; reframed the prefetch horizon as a runtime control variable.
    \item Designed a closed-loop controller with EMA-smoothed stall/eviction signals, a CPU-side delta-adjusted expert forecaster, and a residency-aware two-tier cache with cache-aware token routing.
    \item Raised expert prediction accuracy by 21.8\% and reduced cache-miss latency by 13$\times$ (geometric mean) across 3 MoE models on 3 accelerators. \textit{Under review at ICCAD 2026.}
\end{itemize}

\vspace{0.5ex}
\textbf{MCaM: Multi-Tier KV-Cache Reuse for LLM Inference} \hfill \textit{Nov 2023 -- Aug 2024}
\begin{itemize}[leftmargin=1.2em, itemsep=0.2em, topsep=0.2em]
    \item Built a multi-tier KV-cache system with prefix hashing, deduplication, and memory-aware GPU--CPU offloading.
    \item Reduced redundant prefill computation across requests with similar prefixes under bounded GPU memory.
    \item Achieved 69\% lower prefill latency and 3.3$\times$ higher throughput in prototype evaluation. \textit{Published at ICDCS 2025.}
\end{itemize}

\end{rSection}

% ===================== EXPERIENCE =====================
\begin{rSection}{Industry Experience}\sectionvspace

\textbf{Dcar, Inc. (ByteDance automotive platform)} \hfill \textit{Beijing, China}\\
\textit{ML Systems Engineer Intern} \hfill \textit{Apr 2025 -- Oct 2025}
\begin{itemize}[leftmargin=1.2em, itemsep=0.2em, topsep=0.2em]
    \item Optimized the LLM serving system for Dcar's AI Car Selection service, supporting 400K+ daily requests.
    \item Designed a resource-aware solution combining quantization and runtime offload/prefetch to address cross-node GPU inference inefficiency and oversized model memory footprints on a single 8$\times$H20 node.
\end{itemize}

\vspace{0.5ex}
\textbf{Baidu, Inc. (China's leading search engine company)} \hfill \textit{Beijing, China}\\
\textit{Senior Software Engineer} \hfill \textit{Jul 2020 -- Aug 2024}
\begin{itemize}[leftmargin=1.2em, itemsep=0.2em, topsep=0.2em]
    \item Built and optimized Search Push \& Recommendation (600M pushes/day), DeepQA Search (150M daily visits), and real-time streaming indexing pipelines; promoted twice (T3$\rightarrow$T5) in 3 years.
    \item Reduced end-to-end data availability latency from 24 hours to 30 minutes and incremental indexing latency from 48 hours to 5 minutes via stream+batch pipeline redesigns.
    \item Reduced long-tail service latency by 63\% through Go service decomposition and parallel retrieval.
\end{itemize}

\end{rSection}

% ===================== SKILLS =====================
\begin{rSection}{Technical Skills}\sectionvspace
\begin{tabular}{@{} >{\bfseries}l @{\hspace{2ex}} p{0.84\linewidth} @{}}
Languages &
Python, Go, C++ \\

ML / Serving &
PyTorch, Transformers, vLLM, SGLang, CUDA, KV-cache reuse, deterministic inference \\

Systems / Backend &
Linux, gRPC, FastAPI, Gin, Redis, MySQL, Docker, Git \\

Data / ANN &
Faiss, HNSW, DiskANN, IVF, OOD graph repair, Kafka, Hadoop, Spark
\end{tabular}
\end{rSection}

% ===================== AWARDS =====================
\begin{rSection}{Awards}\sectionvspace
\item Predoctoral Fellowship, University of Connecticut, 2025
\item Baidu Pride Special Award, 2022
\item Baidu Breakthrough Innovation Award, 2021--2022
\item National Scholarship, Hefei University of Technology, 2018, 2019
\end{rSection}

\end{document}
